\documentclass[10pt,a4paper]{article}
\usepackage[utf8]{inputenc}
\usepackage[T1]{fontenc}
\usepackage{lmodern}
\usepackage[ngerman]{babel}
\usepackage[a4paper,margin=9mm]{geometry}
\usepackage{amsmath,amssymb,mathtools}
\usepackage[expansion=false]{microtype}
\usepackage{xcolor}
\usepackage{enumitem}
\pagestyle{empty}
\setlength{\parindent}{0pt}
\setlength{\parskip}{1.4pt}
\setlength{\fboxsep}{2mm}
\setlength{\fboxrule}{0.45pt}
\setlist{nosep,leftmargin=*,topsep=1pt}
\definecolor{ink}{HTML}{173548}
\definecolor{line}{HTML}{78909C}
\newcommand{\R}{\mathbb R}
\newcommand{\norm}[1]{\lVert#1\rVert}
\newcommand{\rang}{\operatorname{rang}}
\newcommand{\vct}[1]{\begin{pmatrix}#1\end{pmatrix}}
\newcommand{\merk}[1]{\par\textbf{Kontrolle:} #1}
\newlength{\tilewidth}
\setlength{\tilewidth}{\dimexpr(\textwidth-3mm)/2-2\fboxsep-2\fboxrule\relax}
\newsavebox{\tilecontent}
\newif\ifkachelueberlauf
\AtEndDocument{\ifkachelueberlauf
  \PackageError{LALG}{Mindestens eine Kachel zu hoch}{Inhalt der gemeldeten Kachel kuerzen.}%
\fi}
\newcommand{\tile}[3]{%
  \sbox{\tilecontent}{\begin{minipage}[t]{\tilewidth}
    \fontsize{9.5}{11}\selectfont\raggedright
    \setlength{\abovedisplayskip}{2.4pt}
    \setlength{\belowdisplayskip}{2.4pt}
    \setlength{\abovedisplayshortskip}{1.4pt}
    \setlength{\belowdisplayshortskip}{1.4pt}
    {\color{ink}\fontsize{11}{12.5}\selectfont\bfseries #1\enspace #2}\par
    \vspace{0.8mm}{\color{line}\hrule}\vspace{1mm}
    #3
  \end{minipage}}%
  \typeout{TILE #1: \the\dimexpr\ht\tilecontent+\dp\tilecontent\relax}%
  \ifdim\dimexpr\ht\tilecontent+\dp\tilecontent\relax>81mm
    \global\kachelueberlauftrue
    \PackageWarning{LALG}{Kachel #1 zu hoch}%
  \fi
  {\color{line}\fbox{\color{black}\begin{minipage}[t][81mm][t]{\tilewidth}
    \usebox{\tilecontent}
  \end{minipage}}}%
}
\newcommand{\rowgap}{\par\vspace{2mm}\noindent}
\newcommand{\head}[2]{%
  {\color{ink}\Large\bfseries LALG\enspace Formelsammlung}\hfill
  {\small\bfseries #1\,/\,2}\par
  {\small #2}\par\vspace{2mm}
}
% Grundlagen: Arbeitsmappe LALG1, Kap. 1--3 und 5--7, Serien 1, 3--5.
% Aufgabenschwerpunkte: Test 1 (13.10.2025). Kapitel 4 bewusst ausgeschlossen.
% Nur Kernformeln und Rechenschritte; keine Beispiele.
% Notation: Skript Kap. 3, 5.2, 7.1--7.4; Test 1; Serien 1, 3--5.
% Vektorpfeile, Skalarprodukt odot, Projektion f, Lot h, Fuss F; kein t_0.
\begin{document}
\head{1}{Notation wie im Skript: Vektoren mit Pfeil, Skalarprodukt $\odot$}
\tile{1}{Geraden}{
\setlength{\belowdisplayskip}{1.4pt}
\setlength{\belowdisplayshortskip}{0.4pt}
\textbf{Parameterform:}
\[
 \boxed{g:\vec x=\vec A+\lambda\vec v},\quad\lambda\in\R.
\]
$\vec A$: Aufpunkt; $\vec v\ne\vec0$: Richtungsvektor.
$\lambda$ (auch $s,t$): frei wählbare Zahl.
Durch $A,B$: $\vec v=\vec B-\vec A$.

\textbf{Koordinatenform in $\R^2$:}
\[
 \vec n=\vct{v_y\\-v_x},\qquad d=-\vec A\odot\vec n,
\]
\[
 \boxed{n_xx+n_yy+d=0}.
\]
$\vec n$: Normalenvektor (senkrecht zur Geraden).
Zurück: Aufpunkt einsetzen und bestimmen.
\[
\vec v=(n_y,-n_x)^T
=\begin{pmatrix}n_y\\-n_x\end{pmatrix}
\]

\textbf{Funktionsform:} $y=mx+d$.

\textbf{Punktprobe:} Punkt einsetzen.\par
\textbf{Schnittpunkt:} beide Gleichungen gemeinsam lösen.
}
\hfill
\tile{2}{Lineare Abhängigkeit}{
\textbf{Für drei Vektoren in $\R^3$:}
\begin{enumerate}
\item \textbf{Matrix aufstellen:} Vektoren als Zeilen.
\[
 \begin{array}{c|ccc}
 \mathrm{GL1}&a_1&a_2&a_3\\
 \mathrm{GL2}&b_1&b_2&b_3\\
 \mathrm{GL3}&c_1&c_2&c_3
 \end{array}
\]
\item \textbf{Position 1 in GL2 auf $0$ bringen:}
$\mathrm{GL2}\leftarrow\mathrm{GL2}-\text{Faktor}\cdot\mathrm{GL1}$.
\item \textbf{Position 1 in GL3 auf $0$ bringen:}
$\mathrm{GL3}\leftarrow\mathrm{GL3}-\text{Faktor}\cdot\mathrm{GL1}$.
\item \textbf{Position 2 in GL3 auf $0$ bringen:}
$\mathrm{GL3}\leftarrow\mathrm{GL3}-\text{Faktor}\cdot\mathrm{GL2}$.
\item \textbf{Nullzeile $(0\ 0\ 0)$: linear abhängig.}
Drei Pivots $\ne0$: linear unabhängig.
\end{enumerate}
\textbf{Faktor:} zu löschender Eintrag geteilt durch Eintrag
der benutzten Zeile in derselben Spalte.

\textbf{Pivot $0$:} zuerst Zeilen tauschen. Kein passender Pivot
möglich: abhängig. Benutzte Zeile unverändert lassen.

\textbf{Linearkombination:} Zeilenoperationen mitführen;
Nullzeile auf die ursprünglichen Vektoren zurückführen.
}
\rowgap
\tile{3}{Dreieck}{
\textbf{Seite $c=AB$ liegt gegenüber $C$.}
\[
 \boxed{\vec M_c=\frac{\vec A+\vec B}{2}}\quad\text{Mittelpunkt},
\]
\[
 \boxed{\vec S=\frac{\vec A+\vec B+\vec C}{3}}\quad\text{Schwerpunkt}.
\]
\textbf{Normalenvektor zur Seite $AB$ in $\R^2$:}
\[
 \vec v=\overrightarrow{AB}=\vec B-\vec A,\qquad
 \vec n=\vct{v_y\\-v_x}.
\]
\textbf{Projektion (Höhenfusspunkt):}
\[
 \boxed{\vec C'=\vec F_c=\vec C-
 \frac{(\vec C-\vec A)\odot\vec n}{\vec n\odot\vec n}\,\vec n}.
\]

\textbf{Höhenvektor und Höhe:}
$\vec h_c=\vec C-\vec F_c$, $\boxed{h_c=\norm{\vec h_c}}$.

\textbf{Fläche:} $\mathcal A=\tfrac12\norm{\overrightarrow{AB}}\,h_c$.
}
\hfill
\tile{4}{Abstand und Projektion auf Gerade}{
\textbf{Gegeben:} $g:\vec x=\vec A+\lambda\vec v$, Punkt $\vec P$.
$\vec A$: Aufpunkt; $\vec v\ne\vec0$: Richtung. Gilt in $\R^2,\R^3$.

\textbf{Verbindungsvektor:} $\overrightarrow{AP}=\vec P-\vec A$.

\textbf{Projektion in Geradenrichtung:}
\[
 \boxed{\vec f=\frac{\overrightarrow{AP}\odot\vec v}
 {\vec v\odot\vec v}\,\vec v}.
\]
\textbf{Projektion:} $\boxed{\vec P'=\vec F=\vec A+\vec f}$.

\textbf{Spiegelung:} $\boxed{\vec P''=2\vec P'-\vec P}$.

\textbf{Lotvektor:} $\vec h=\vec P-\vec P'$.
\textbf{Abstand:} $\boxed{h=\norm{\vec h}}$.

\textbf{Vektorprojektion von $\vec b$ auf $\vec a\ne\vec0$:}
\[
 \vec f=\frac{\vec b\odot\vec a}{\vec a\odot\vec a}\,\vec a,
 \qquad\vec h=\vec b-\vec f.
\]
$\vec f$: parallel zu $\vec a$; $\vec h$: senkrecht zu $\vec a$.
}
\rowgap
\tile{5}{Vektoren, Länge und Winkel}{
\textbf{Verbindung:} $\overrightarrow{AB}=\vec B-\vec A$.
\textbf{Abstand der Punkte:} $\norm{\vec B-\vec A}$.

\textbf{Skalarprodukt} (Orthonormalbasis):
\[
 \boxed{\vec a\odot\vec b=a_1b_1+a_2b_2+a_3b_3}.
\]
\textbf{Norm / Länge:}
\[
 \boxed{\norm{\vec a}=\sqrt{a_1^2+a_2^2+a_3^2}}.
\]
\textbf{Normieren:} $\vec a/\norm{\vec a}$ für $\vec a\ne\vec0$.

\textbf{Senkrecht:} $\vec a\odot\vec b=0$.

\textbf{Zwischenwinkel $\varphi$} ($\vec a,\vec b\ne\vec0$):
\[
 \boxed{\varphi=\arccos\!\left(\frac{\vec a\odot\vec b}
 {\norm{\vec a}\,\norm{\vec b}}\right)}.
\]
Kleiner Geradenwinkel: Betrag im Zähler.
Für Winkel in Grad: Taschenrechner auf \textbf{DEG}.
}
\hfill
\tile{12}{Vektorraum prüfen}{
Für eine Menge $V\subseteq\R^n$ mit den üblichen Rechenregeln:
\begin{enumerate}
\item \textbf{Nullvektor:} $\vec0$ in die Bedingung einsetzen.
Ist $\vec0\notin V$: \textbf{kein Vektorraum}.
\item \textbf{Addition:} zwei beliebige Vektoren
$\vec u,\vec v\in V$ wählen. Prüfen, ob ihre Summe die Bedingung von $V$ erfüllt:
\[
 \boxed{\vec u+\vec v\in V}.
\]
\item \textbf{Multiplikation:} beliebiges $\lambda\in\R$
und $\vec u\in V$ wählen. Prüfen, ob das Vielfache die Bedingung von $V$ erfüllt:
\[
 \boxed{\lambda\vec u\in V}.
\]
\end{enumerate}
\textbf{Alle drei erfüllt: Vektorraum.}
Eine Bedingung verletzt: kein Vektorraum.

\textbf{Nachweis:} für beliebige Vektoren und Zahlen zeigen.
\textbf{Widerlegung:} ein Gegenbeispiel reicht.

\textbf{Bei Parameterdarstellung:} Summe und Vielfaches müssen
wieder dieselbe Form mit zulässigem Parameter haben.
Bei beschränktem Parameter den erlaubten Bereich prüfen.
Verschobene Darstellungen zuerst vereinfachen.
}
\newpage
\head{2}{Gauss, Vektorprodukt, Ebenen und LGS-Lösungen}
\tile{6}{Gauss und Lösbarkeit}{
\textbf{Wie in Serie 5:} $A$: Koeffizientenmatrix;
$A_e$: erweiterte Matrix (mit rechter Seite).
$L_1,L_2,\ldots$: Gleichungen bzw.\ Zeilen.

\textbf{Erlaubt:} Zeilen tauschen, mit Zahl $\ne0$ multiplizieren,
Vielfaches einer anderen Zeile addieren.
\[
 L_i\leftarrow L_i+\lambda L_j\qquad(i\ne j).
\]
\textbf{Vorgehen:}
\begin{enumerate}
\item Pivot (erster Eintrag $\ne0$ einer Zeile) wählen.
Falls nötig, Zeilen tauschen.
\item Darunter Nullen erzeugen; benutzte Zeile beibehalten.
\item Nächste Zeile; am Ende von unten rückwärts einsetzen.
\end{enumerate}
\textbf{Rang:} Anzahl Pivots. Bei $n$ Unbekannten:
\begin{itemize}
\item $0=c$, $c\ne0$: \textbf{keine Lösung}.
\item Kein Widerspruch, $\rang(A)=n$: \textbf{eine Lösung}.
\item Kein Widerspruch, $\rang(A)<n$: \textbf{unendlich viele};
$n-\rang(A)$ freie Variablen.
\end{itemize}
$0=0$ streichen. Homogen (rechte Seite $0$): immer lösbar.
}
\hfill
\tile{7}{Vektorprodukt, Fläche und Volumen}{
\textbf{Kreuzprodukt:}
\[
 \boxed{\vec a\times\vec b=\vct{
 a_2b_3-a_3b_2\\a_3b_1-a_1b_3\\a_1b_2-a_2b_1}}.
\]
Ergebnis steht senkrecht auf $\vec a$ und $\vec b$.
Vertauschen ändert das Vorzeichen.

\textbf{Flächen:}
\[
 \mathcal A_{\text{Parallelogramm}}=\norm{\vec a\times\vec b},
\]
\[
 \boxed{\mathcal A_{\triangle ABC}
 =\tfrac12\norm{\overrightarrow{AB}\times\overrightarrow{AC}}}.
\]
\textbf{Spatvolumen:}
\[
 \boxed{V=|\vec a\odot(\vec b\times\vec c)|}.
\]
\textbf{Abhängigkeit:} $\vec a\times\vec b=\vec0$
 genau dann, wenn $\vec a,\vec b$ abhängig sind.
Drei Vektoren in $\R^3$ komplanar genau dann,
wenn $\vec a\odot(\vec b\times\vec c)=0$.
}
\rowgap
\tile{8}{Ebenen darstellen und umwandeln}{
\textbf{Durch drei nicht kollineare Punkte:}
$\vec u=\vec B-\vec A$, $\vec v=\vec C-\vec A$.
\[
 \boxed{E:\vec x=\vec A+\lambda\vec u+\nu\vec v}.
\]
\textbf{Parameter $\to$ Koordinaten durch Einsetzen:}
\begin{enumerate}[label=\arabic*.]
\item Komponentengleichungen aufschreiben.
\item Nach $\lambda$ (oder $t$) auflösen und in die anderen Gleichungen einsetzen; auch $\nu$ beseitigen.
\item Vereinfachen bis $n_1x+n_2y+n_3z+d=0$.
\end{enumerate}

\textbf{Parameter $\to$ Normalenform:}
\[
 \vec n=\vec u\times\vec v\ne\vec0,\qquad
 \boxed{(\vec x-\vec A)\odot\vec n=0}.
\]
\textbf{Koordinatenform wie im Skript:}
\[
 \boxed{n_1x+n_2y+n_3z+d=0},
 \qquad d=-\vec A\odot\vec n.
\]

\textbf{Koordinaten $\to$ Parameter:} Eine Variable isolieren,
zwei übrige als Parameter setzen; nach Parametern sammeln.

\textbf{Punktprobe:} Punkt in Koordinatenform einsetzen.
}
\hfill
\tile{9}{Abstand und Projektion auf Ebene}{
\textbf{Gegeben:} Aufpunkt $\vec A$, Normalenvektor $\vec n\ne\vec0$,
Punkt $\vec P$. Bei Parameterform $\vec n=\vec u\times\vec v$.

\textbf{Hesse-Normalenform wie in Kap. 7:}
\[
 \boxed{h(\vec P)=\frac{(\vec P-\vec A)\odot\vec n}{\norm{\vec n}}}.
\]
$h(\vec P)$ kann negativ sein. \textbf{Abstand:} $|h(\vec P)|$.

Bei $ax+by+cz+d=0$, $\vec P=(p_1,p_2,p_3)^T$:
\[
 h(\vec P)=\frac{ap_1+bp_2+cp_3+d}{\sqrt{a^2+b^2+c^2}}.
\]
\textbf{Projektion (Lotfusspunkt):}
\[
 \boxed{\vec P'=\vec F=\vec P-h(\vec P)\,\frac{\vec n}{\norm{\vec n}}}.
\]
\textbf{Spiegelung:} $\boxed{\vec P''=2\vec P'-\vec P}$.

\textbf{Ursprungsabstand:} $|d|/\sqrt{a^2+b^2+c^2}$.
}
\rowgap
\tile{10}{Schnittmengen mit Ebenen}{
\textbf{Gerade mit Ebene:}
$\vec x=\vec A+\lambda\vec v$ in $\vec n\odot\vec x+d=0$ einsetzen:
\[
 (\vec n\odot\vec v)\lambda+\vec n\odot\vec A+d=0.
\]
\begin{itemize}
\item $\vec n\odot\vec v\ne0$: $\lambda$ berechnen,
$\vec S=\vec A+\lambda\vec v$.
\item $\vec n\odot\vec v=0$: Aufpunkt einsetzen.
Erfüllt er die Ebene, liegt die Gerade darin; sonst kein Schnitt.
\end{itemize}
\textbf{Zwei Ebenen $\vec n_i\odot\vec x+d_i=0$:}
\begin{itemize}
\item $\vec n_1\times\vec n_2\ne\vec0$: Schnittgerade.
Beide Gleichungen als LGS lösen; eine freie Variable.
\item $\vec n_2=k\vec n_1$: bei $d_2=kd_1$ identisch,
sonst parallel ohne Schnitt ($k$: reelle Zahl).
\end{itemize}
\textbf{Drei Ebenen:} Gauss auf alle Gleichungen.
Widerspruch: kein Schnitt. Sonst:
\[
 \begin{array}{c|ccc}
 \rang(A)&3&2&1\\\hline
 \text{Schnitt}&\text{Punkt}&\text{Gerade}&\text{Ebene}
 \end{array}
\]
\textbf{Homogen:} alle $d_i=0$; Ursprung immer im Schnitt.
}
\hfill
\tile{11}{LGS-Lösung als Gerade oder Ebene}{
\begin{enumerate}
\item \textbf{Freie Variablen bestimmen:} Variablen ohne Pivot
als Parameter $\lambda,\nu$ setzen.
\item \textbf{Übrige Variablen berechnen:} von unten rückwärts
einsetzen; nach Parametern sammeln.
\item \textbf{Aufpunkt $\vec P$:} alle freien Parameter $0$ setzen.
\item \textbf{Richtungsvektoren $\vec u,\vec v$:} rechte Seite
des LGS $0$ setzen. Je einen freien Parameter $1$, übrige $0$;
LGS lösen.
\end{enumerate}
\textbf{Eine freie Variable: Schnittgerade}
\[
 \boxed{\vec x=\vec P+\lambda\vec u},\quad\lambda\in\R.
\]
\textbf{Zwei freie Variablen: gemeinsame Ebene}
\[
 \boxed{\vec x=\vec P+\lambda\vec u+\nu\vec v},
 \quad\lambda,\nu\in\R.
\]
$\vec P$: Aufpunkt; $\vec u,\vec v$: Richtungsvektoren.
}
\end{document}
